% ----------------------------------------------------------------------
% Capítulo: Introdução
%
% A introdução é a primeira parte da estrutura textual e deve iniciar na
% página 1, com numeração em algarismos arábicos. Usar o comando
% \chapter para produzir o título numerado e siga as diretrizes de
% espaçamento e formatação. Os títulos de seções subsequentes devem
% respeitar a numeração progressiva.

\chapter{Introdução}

O estudo do tempo até o parto constitui um problema relevante na pesquisa em saúde 
perinatal, especialmente quando diferentes desfechos podem ocorrer ao longo da mesma 
escala temporal. Entre esses desfechos destacam-se o óbito fetal e o nascimento vivo, 
eventos terminais e mutuamente exclusivos que caracterizam a evolução temporal da 
gestação. Na epidemiologia aplicada, esses desfechos são frequentemente analisados 
em modelos separados, o que enfraquece a interpretação temporal e pode gerar 
inconsistências quando um evento impede a ocorrência do outro. O tempo gestacional, 
portanto, se ajusta a uma estrutura de riscos competitivos, na qual riscos 
causa-específicos e funções de incidência acumulada (CIF, do inglês 
\textit{Cumulative Incidence Function}) são mais apropriados do que métodos 
de sobrevivência de evento único \cite{satagopan2004note,finegray1999,
naimi2016cumulative,smith2016competing}.

Nesta dissertação, o interesse central é avaliar o tempo até o parto em
gestações de residentes no estado de São Paulo, considerando o óbito fetal e o
nascimento vivo como eventos competitivos. Essa abordagem permite tratar o nascimento 
vivo não apenas como um registro de encerramento da gestação, mas como um desfecho 
que impede a observação posterior do óbito fetal naquela mesma gestação. Ao mesmo 
tempo, a inclusão dos nascidos vivos aproxima o estudo de outro tema importante 
em obstetrícia, sendo este, o problema da prematuridade. Dessa forma, o trabalho 
não se restringe à mortalidade fetal, mas busca compreender a dinâmica temporal 
do parto em presença de desfechos competitivos. O estado de São Paulo, por sua vez, 
apresenta um cenário estratégico para análise devido à sua população numerosa e 
heterogênea, com grande volume de nascimentos, ampla rede hospitalar, desigualdades 
regionais, diferentes níveis de organização da rede assistencial, aspectos que 
podem ser parcialmente investigados (no âmbito deste estudo) por meio de dados 
públicos do Sistema de Informações sobre Nascidos Vivos (SINASC) e do Sistema de 
Informações sobre Mortalidade (SIM).

Além disso, é necessária atenção à definição dos casos, dado que diferentes critérios 
podem ser utilizados na caracterização dos desfechos, principalmente em relação ao 
óbito fetal, cujas definições podem variar entre sistemas nacionais e critérios 
internacionais de comparabilidade \cite{Andrade2005,Lawn2016}.

\section{Definição dos desfechos perinatais e cenário global}

No Brasil, a Declaração de Óbito deve ser emitida para óbito fetal quando a
gestação tiver duração igual ou superior a 20 semanas, ou quando o feto
apresentar peso corporal igual ou superior a $500\,\text{g}$, ou estatura igual
ou superior a $25\,\text{cm}$ \cite{ministerioSaudeDO2022}. Este é o critério
que orienta a definição adotada neste estudo e também é compatível com a
normatização nacional para emissão da Declaração de Óbito \cite{Andrade2005}. O 
nascimento vivo, por sua vez, é definido como a expulsão ou extração completa 
do produto da concepção que, após a separação do corpo materno, apresenta 
respiração ou qualquer outro sinal de vida, independentemente da duração da 
gestação \cite{ministerioSaudeDNV2022}. Portanto, mesmo nas idades gestacionais 
mais precoces, um recém-nascido com sinais de vida é registrado como nascido 
vivo. Essas definições são importantes para compreensão dos desfechos perinatais 
ao longo da gestação e consistência das informações apresentadas.

Em ambas definições, a idade gestacional é de extrema importância, pois a interpretação 
clínica e epidemiológica dos desfechos depende do momento gestacional em que o parto 
ocorre. A prematuridade entra nesse raciocínio por estar contida no conjunto dos
nascidos vivos. A Organização Mundial da Saúde (OMS) define nascimento prematuro
como aquele ocorrido antes de 37 semanas completas de gestação, tendo como subdivisões 
a prematuridade extrema, muito prematuro e prematuridade moderada a tardia
\cite{WHO2023Preterm}. Estimativas globais de 2020 indicaram a ocorrência de 13,4 
milhões de nascimentos prematuros, com importante variação entre países e
regiões \cite{Ohuma2023PretermEstimates}. O relatório \textit{Born Too Soon},
publicado pela OMS, Fundo das Nações Unidas para a Infância (UNICEF) e
parceiros, reforçou que a carga da prematuridade permanece elevada, com pouco
avanço global na última década e forte relação com desigualdades no acesso e na
qualidade do cuidado materno e neonatal \cite{BornTooSoon2023}. De forma semelhante, 
estimativas internacionais recentes mostram que a carga de óbito fetal permanece 
elevada e distribuída de modo desigual entre países e sistemas de saúde \cite{Lawn2016,UnitedNationsChildMortality,Neglectedt}. 
Esses elementos reforçam a importância de analisar não apenas a ocorrência dos desfechos, 
mas também o momento gestacional em que eles acontecem e o contexto assistencial 
em que a gestação foi acompanhada.

Diante da elevada carga global de prematuridade e óbito fetal, bem como das 
desigualdades observadas, torna-se importante compreender o panorama brasileiro 
considerando os desfechos perinatais em questão.

\section{Panorama brasileiro}

No Brasil, a mortalidade fetal e a prematuridade seguem sendo importantes problemas 
de saúde perinatal, fortemente relacionados às condições de saúde materna, à qualidade 
do cuidado pré-natal e à assistência ao parto. Em estudo nacional recente com dados do 
Sistema de Informações sobre Mortalidade (SIM), considerando óbitos fetais com idade 
gestacional igual ou superior a 20 semanas, Rocha et al. observaram que esses óbitos 
representaram 1,14\% dos nascimentos ocorridos no período de 1996 a 2021, com taxa 
global (considerando todo o período) de 11,4 óbitos fetais por 1.000 nascimentos 
\cite{rocha2025space}. Apesar de Rocha et al. observarem uma redução de aproximadamente 
20\% na taxa de mortalidade fetal nesse período, a queda não foi homogênea no território 
nacional, sendo observadas maiores taxas e reduções mais lentas nas regiões Norte, 
Nordeste e Centro-Oeste, evidenciando desigualdades regionais \cite{rocha2025space}. 
Além disso, cerca de 93\% dos óbitos ocorreram antes do parto, aspecto que reforça a 
relevância da vigilância durante a gestação.

Esse padrão de desigualdade também aparece quando se observa o nascimento vivo pela 
perspectiva da prematuridade. Análises nacionais baseadas no Sistema de Informações 
sobre Nascidos Vivos (SINASC), envolvendo mais de 25,5 milhões de nascidos vivos 
entre 2014 e 2023, indicaram aumento da prevalência de prematuridade de 11,3\% em 
2014 para 11,9\% em 2023 \cite{Victor2025PretermBrazil}. As maiores prevalências 
foram observadas entre mães com 35 anos ou mais, mães com menos de 20 anos, mulheres 
sem escolaridade formal e nas regiões Norte e Nordeste \cite{Victor2025PretermBrazil}. 
Esse padrão é coerente com o Boletim Epidemiológico do Ministério da Saúde sobre 
nascimentos prematuros no período de 2012 a 2022, que descreve a prematuridade 
como evento associado a fatores sociodemográficos, condições maternas, cuidado 
pré-natal e características da gestação \cite{BrasilBoletimPrematuridade2024}. 
Além disso, o Brasil registrou 2.537.511 nascidos vivos em 2023, dos quais 304.501 foram 
prematuros, correspondendo a 12\% do total \cite{OOBrPainelVigilancia}. 
Esses números ajudam a dimensionar a magnitude da prematuridade no país e reforçam 
a relevância epidemiológica dos desfechos perinatais no contexto brasileiro.

Essas evidências indicam que os desfechos perinatais em questão apresentam importante 
variabilidade em território nacional, mostrando a importância de análises que considerem 
diferenças regionais e contextos assistenciais específicos. Assim, o estado de São Paulo 
apresenta um cenário de grande interesse de investigação dado o elevado volume de 
nascimentos e à complexidade de sua organização regional de saúde. 

\section{Contexto do estado de São Paulo e das RRAS}

O recorte do estado de São Paulo, embora sendo parte desse cenário nacional,
acrescenta uma dimensão própria. O estado reúne grande volume de nascimentos e 
uma rede assistencial complexa, com diferenças internas de organização, capacidade 
instalada e fluxos regionais de atenção. A regionalização da saúde no estado é 
organizada, dentre outras, por Regiões de Saúde e Redes Regionais de Atenção à 
Saúde (RRAS), estruturas voltadas à integração de ações e serviços no território 
\cite{SESSPRegionalizacao2026}. Para a análise de dados de 2023, essa organização 
regional é importante porque permite avaliar se o risco de óbito fetal e a distribuição 
do nascimento vivo variam entre agregados territoriais que compartilham redes de 
referência e condições assistenciais semelhantes.

Em 2023 o estado de São paulo registrou 503.903 nascidos vivos, do quais 60.468 foram 
prematuros, o que corresponde a aproximadamente 12\% do total \cite{OOBrPainelVigilancia}. 
Esses valores situam o estado dentro do cenário nacional de prematuridade e reforçam 
a importância de estudar o tempo até o parto em uma população numerosa, em linha 
com estudos nacionais que descrevem a prematuridade como evento frequente e 
socialmente distribuído no Brasil \cite{Victor2025PretermBrazil,BrasilBoletimPrematuridade2024}. 
Entretanto, para os óbitos fetais, a contextualização estadual exige cautela adicional, 
pois as estimativas disponíveis podem variar conforme a definição adotada para o 
desfecho, especialmente em relação à idade gestacional mínima considerada. Dados 
disponíveis e estudos recentes conduzidos nesse recorte adotaram definições mais 
restritivas, como óbitos a partir de 22 semanas de gestação. Assim, a magnitude 
estadual do óbito fetal será apresentada diretamente nas análises do estudo, de 
acordo com a definição adotada para a construção da população analisada. 

A necessidade de considerar heterogeneidade territorial também é sustentada pela
literatura. No município de São Paulo, um estudo identificou diferenças
intraurbanas de mortalidade fetal segundo agrupamentos de vulnerabilidade
social, indicando que a distribuição dos óbitos não se resume ao volume
populacional das áreas analisadas \cite{Marques2021FetalSocialVulnerabilitySP}.
Além disso, o estudo FetRiskS, realizado em hospitais públicos do município de
São Paulo com dados coletados entre 2019 e 2023, descreveu o óbito fetal como
evento relacionado a fatores socioeconômicos, acesso inadequado ao pré-natal,
condições maternas, alterações placentárias, restrição de crescimento fetal e
malformações congênitas \cite{Marques2026FetRiskS}. Embora esses achados não possam 
ser refletidos para o conjunto estadual e das RRAS, eles apontam para a importância 
de considerar que a região de residência possa aproximar condições contextuais 
compartilhadas relacionadas à organização da rede, ao acesso aos serviços e às 
condições de assistência durante a gestação e o parto.

Assim, o estado de São Paulo não foi escolhido apenas pelo tamanho da base de 
dados, mas também por reunir regiões com diferentes formas de organização da 
rede de saúde, capacidade assitencial e condições de acesso ao cuidado. Este cenário 
é adequado para investigar se o risco e a dinâmica temporal do óbito fetal e do nascimento 
vivo varia entre regiões, mesmo após considerar as características maternas, fetais, 
obstétricas e assistenciais disponíveis no SIM-DOFET e no SINASC. Caso a região 
de residência permaneça associada aos riscos estimados após o ajuste, esse resultado 
não deve ser interpretado como prova de que a região, por si só, determine o 
desfecho. De forma mais cautelosa, essa associação deve ser vista como possível 
indicativo de diferenças contextuais não captadas pelas variáveis individuais 
disponíveis, reforçando a importância de considerar, em conjunto, fatores maternos, 
fetais, obstétricos, assistenciais e territoriais associados aos desfechos 
gestacionais.

\section{Fatores associados aos desfechos gestacionais}

A interpretação dos desfechos gestacionais requer considerar o caráter multifatorial 
desses eventos. Tanto o óbito fetal quanto a prematuridade resultam da interação 
entre fatores maternos, fetais, placentários, obstétricos, assistenciais e sociais 
que atuam de forma simultânea e, muitas vezes, dependente ao longo da gestação. 
Documentos clínicos e revisões da literatura apontam de forma recorrente a associação 
de óbito fetal com condições como hipertensão crônica, pré-eclâmpsia, diabetes, 
obesidade, tabagismo, idade materna avançada, gestação múltipla, histórico obstétrico 
desfavorável e vulnerabilidade social \cite{Lawn2016,ACOG2020Stillbirth,Sun2018,zile2019maternal}. 
De maneira similar, a prematuridade compartilha parte importante desses fatores, embora 
seja registrada entre os nascidos vivos. Estudos observacionais e revisões da literatura 
descrevem o nascimento prematuro como um desfecho associado a diferentes condições 
maternas, históricas, obstétricas, ambientais e socioeconômicas, frequentemente 
relacionadas ao longo da gestação \cite{Mitrogiannis2023PretermRisk}. No Brasil, 
evidências recentes mostram maior frequência de prematuridade nos extremos de idade 
materna, em contextos de menor escolaridade, em gestações de maior complexidade e em 
grupos com menor acesso a cuidado pré-natal adequado 
\cite{Victor2025PretermBrazil,BrasilBoletimPrematuridade2024}. Assim, fatores como 
baixa escolaridade, vulnerabilidade socioeconômica e acesso limitado ao pré-natal 
não devem ser interpretados apenas como características individuais, mas também 
como marcadores de trajetórias sociais e assistenciais distintas durante a gestação.

Entre os fatores fetais e placentários, destacam-se restrição de crescimento
fetal, malformações congênitas, anomalias genéticas, insuficiência placentária e
descolamento prematuro de placenta \cite{gardosi2013maternal,smith2007stillbirth,pinar2014placental}. 
A restrição de crescimento fetal é particularmente importante porque pode permanecer 
não identificada durante o pré-natal e se manifestar apenas no desfecho da gestação. 
Além disso, muitos desses fatores atuam de forma dependente. A obesidade e a diabetes, 
por exemplo, podem aumentar o risco de distúrbios hipertensivos, enquanto distúrbios 
hipertensivos podem comprometer a função placentária e favorecer insuficiência 
placentária que, por sua vez, pode antecipar o parto ou resultar em óbito fetal. Essa 
dependência reforça a importância de analisar simultaneamente o óbito fetal e o 
nascimento vivo ao longo da gestação, visto que certas características podem estar 
associadas tanto ao risco de óbito intrauterino quanto ao momento gestacional em 
que o parto ocorre.

As variáveis disponíveis no SIM-DOFET e no SINASC permitem representar apenas parte
desse conjunto de fatores descritos na literatura. Idade materna, escolaridade, 
raça/cor, tipo de gestação, histórico reprodutivo, número de filhos vivos 
e mortos previamente, consultas de pré-natal e momento de início do pré-natal 
representam características maternas, sociais, obstétricas e assistenciais frequentemente 
associadas aos desfechos gestacionais \cite{Victor2025PretermBrazil,Marques2026FetRiskS,
ACOG2020Stillbirth,Silva2026PrenatalInequities}. Da mesma forma, peso ao nascer ou 
peso fetal, sexo, presença de anomalia e Apgar no quinto minuto, quando disponível, 
representam características fetais e neonatais relevantes, embora algumas delas 
sejam medidas no momento do desfecho e, portanto, devam ser interpretadas como 
marcadores do processo gestacional ou do estado ao nascimento, e não necessariamente 
como causas antecedentes. De maneira similar, variáveis como tipo de parto, local 
de ocorrência, indução do trabalho de parto e cesariana antes do início do trabalho 
de parto refletem decisões e condições assistenciais próximas ao encerramento da 
gestação, o que exige cautela na interpretação. 

Embora o SIM-DOFET e o SINASC registrem informações essenciais sobre a mãe, a 
gestação, o feto, o nascimento e o óbito, essas bases representam apenas parte 
dos processos clínicos, assistenciais e sociais envolvidos nos desfechos 
gestacionais. O cuidado pré-natal é um exemplo dessa limitação. Embora seja uma 
fator importante da assistência durante a gestação, sua interpretação no 
SINASC exige cautela, pois a base informa principalmente o número de consultas 
registradas, e não a qualidade clínica do acompanhamento prestado. Em estudo 
nacional com dados de 2023, Silva et al. observaram que a realização de ao 
menos uma consulta de pré-natal apresentava cobertura considerável entre os 
nascidos vivos, mas que apenas 67,4\% dos registros indicavam oito ou mais 
consultas \cite{Silva2026PrenatalInequities}. Esse contraste sugere que a 
cobertura inicial do pré-natal é elevada no país, embora o acompanhamento mais 
completo ainda apresente desigualdades por escolaridade materna, raça/cor, 
idade e região.

Além do pré-natal, outros aspectos relevantes do acompanhamento gestacional e da 
assistência ao parto não estão disponíveis de forma completa e padronizada nas 
bases utilizadas, como resultados de exames, condutas clínicas adotadas durante 
a gestação e o parto e condições socioeconômicas individuais mais detalhadas. 
Mesmo variáveis de alta relevância para este estudo, como a idade 
gestacional, podem apresentar limitações de concordância e diferenças regionais 
de mensuração, apesar da elevada cobertura nacional do SINASC 
\cite{Szwarcwald2019SINASC}. Por isso, os resultados devem ser interpretados 
como associações estimadas a partir de registros populacionais, e não como uma 
evidência direta dos mecanismos clínicos responsáveis por cada desfecho. Esse 
conjunto de limitações reforça a importância de abordagens metodológicas capazes 
de considerar simultaneamente múltiplos fatores associados, restrições derivadas 
das bases de dados e possíveis diferenças contextuais entre regiões.

\section{Relevância metodológica}

A análise de sobrevivência é adequada para este problema porque a idade gestacional 
representa o tempo até a ocorrência do desfecho. Entretanto, quando o interesse 
envolve óbito fetal, o nascimento vivo não deve ser tratado como censura. Ao ocorrer 
o nascimento vivo, a gestação se encerra e o óbito fetal deixa de ser possível. 
Desconsiderar essa estrutura pode distorcer a interpretação do risco acumulado ao 
longo da gestação, especialmente em idades gestacionais nas quais os nascimentos 
vivos se tornam cada vez mais frequentes \cite{naimi2016cumulative,smith2016competing}. 
Os modelos de riscos competitivos oferecem uma forma mais coerente de lidar com isso, 
pois reconhecem que diferentes desfechos competem pelo encerramento do tempo de 
seguimento \cite{satagopan2004note,finegray1999,naimi2016cumulative,Austin2016CompetingRisks}. 

Além da presença de riscos competitivos, os dados possuem uma estrutura regional 
formada pelas RRAS, de forma que gestantes residentes em uma mesma RRAS estão 
inseridas em um mesmo contexto territorial e assistencial, podendo compartilhar 
condições semelhantes de acesso aos serviços, disponibilidade de recursos de saúde 
e características socioambientais não mensuradas pelas variáveis disponíveis. Essa 
interpretação é coerente com a organização regional da assistência em saúde e com 
evidências de desigualdades territoriais em desfechos perinatais 
\cite{rocha2025space,SESSPRegionalizacao2026,Marques2021FetalSocialVulnerabilitySP}. 
Nesse contexto, modelos com fragilidade compartilhada permitem incorporar a 
dependência entre observações pertencentes a uma mesma RRAS por meio de efeitos 
aleatórios associados às unidades regionais. Esses efeitos representam fontes de 
heterogeneidade não observada e acomodam a correlação intra-regional 
\cite{clayton1978model,vaupel1979impact,balan2020tutorial}. Quando essa 
heterogeneidade apresenta estrutura territorial, prioris espaciais do tipo ICAR 
permitem que regiões vizinhas compartilhem informação durante a estimação 
\cite{besag1991,rueheld2005}. Assim, após o ajuste pelas covariáveis medidas, a 
presença de variação residual entre RRAS pode indicar que diferenças regionais 
não captadas individualmente ou outros fatores não considerados ainda influenciam 
o risco dos desfechos analisados.

Essa estrutura metodológica é coerente com o principal problema do estudo: analisar 
o tempo até o parto, considerando que óbito fetal e nascimento vivo são eventos 
competitivos e que os dados estão organizados em um território regionalizado. A 
contribuição esperada não é apenas estimar associações individuais, mas oferecer uma 
forma de análise que respeite a natureza dos desfechos perinatais e a organização 
territorial da assistência em saúde. 

Os capítulos seguintes desenvolvem essas ideias de forma mais detalhada. O Capítulo 2 
apresenta os objetivos gerais e específicos do estudo. O Capítulo 3 descreve as 
fontes de dados, o processo de construção das bases, as variáveis utilizadas e a 
fundamentação estatística do modelo. O Capítulo 4 reúne os resultados das análises 
descritivas e dos modelos ajustados.
